\documentclass[10pt,a4paper]{article}
\usepackage[T1]{fontenc}
\usepackage[margin=20mm]{geometry}
\usepackage{lmodern,mathrsfs,amsmath,amssymb,amsthm,mathtools,microtype}
\usepackage[hidelinks]{hyperref}
\newtheorem{theorem}{Theorem}
\title{A closed failure component in the\protect\\Erd\H{o}s--Straus shifted-factor system}
\author{The Clankers}
\date{18 September 2026}
\begin{document}\maketitle
\begin{abstract}
We retain every factor edge and both divisor channels in the universal
external-nonresidue graph. We give an exact finite reachability criterion and
an integral inverse for its cycles. The hard prime 2521 has a closed seven-cycle
whose complete local selectors all fail. Thus a proposed local-cycle rescue
principle is false, although ES holds at that prime. Composite shifts, original
multiplicities and the same-prime repair are explicit. No global ES result or
historical-priority claim is asserted.
\end{abstract}

\paragraph{The original source.}
Let $p$ be prime in $1,121,169,289,361,529\pmod{840}$ and put
\[
 V=\{q<p:q\text{ prime},(q/p)=-1\},\quad
 \sigma_q=\begin{cases}1&q=3\pmod4,\\3&q=1\pmod4,\end{cases}
 \quad R_q=\sigma_q q,\quad A_q=(p+R_q)/4.
\]
Here $V\ne\varnothing$, $p/4<A_q<p$, $\gcd(pA_q,R_q)=1$, and
$(A_q/p)=-1$. A nonresidue integer below $p$ supplies a nonresidue prime;
the hard-prime characters of $2,3,5,7$ are positive. The source statements
follow from $4A_q\equiv\sigma_q q\pmod p$ and $\sigma_q\in\{1,3\}$.
Hence \emph{every} nonresidue prime factor $r\mid A_q$ lies in $V\setminus\{q\}$.
Keep every edge $q\to r$, with weight $v_r(A_q)$; their sum at each vertex is odd.
The basic graph is inherited from \cite[\S\S1--7]{Graph}; selecting just one
factor would lose the complete reachability question.

Write $A_q=\prod_t t^{e_t}$ and retain all $\beta_t\in[-e_t,e_t]$.
The source word and divisor are
$v=\prod t^{\beta_t}\pmod{R_q}$ and $u=\prod t^{e_t+\beta_t}$.
The complete gates are
\[
 E_q=\{u\mid A_q^2:R_q\mid4u+1\},\qquad
 M_q=\{u\mid A_q^2:R_q\mid u+A_q\}.
\]
Their targets are $-p^{-1}$ and $-1$; inversion of the same exponent vector
adds the equivalent exterior target $-p$, not a third channel.
The original normalization is
\[
 d=\gcd(A_q,u),\quad(h,r,s)=(d^2/u,u/d,A_q/d).
\]
A gate hit gives
\[
 E:\ \kappa=(pr+s)/R_q,\quad(A_q,hs\kappa,phr\kappa),
 \qquad
 M:\ \lambda=(r+s)/R_q,\quad(A_q,phs\lambda,phr\lambda).
 \tag{1}
\]
Valuations prove $A_q=hrs$, $u=hr^2$, and coprimality. Units give integral
quotients; clearing reciprocals proves $4/p$. The ordered inverses are
$u=A_q^2/(R_qy-pA_q)$ in E and $u=pA_q^2/(R_qy-pA_q)$ in M.
These are inherited Type I/II coordinates \cite[\S2]{ET}.
Auxiliary $A_q>p/2$ vertices remain: M is then impossible, while a passing E
state returns to the first-half atlas by the exact switch
$(h,r,s,\kappa)\mapsto(h,\kappa,s,r)$, swapping its first two denominators.

\begin{theorem}[Complete closed-failure test]
Let $\mathscr H=\{q:E_q\cup M_q\ne\varnothing\}$, $U_0=V\setminus\mathscr H$,
and $U_{j+1}=\{q\in U_j:N^+(q)\subseteq U_j\}$, using every outgoing factor.
The terminal set $U$ is reached after at most $|V|$ iterations and is the
largest closed set of locally failed vertices. A vertex outside $U$ has a path
to an actual hit of length at most $|V|-1$. $U$ is nonempty exactly when a sink
strongly connected component has only failed vertices.
\end{theorem}
\begin{proof}
Every strict iteration removes a vertex. Every closed failed set remains in
all $U_j$. Removed vertices have a successor removed earlier, furnishing a
strictly decreasing removal rank and a terminating hit path. Conversely,
backwards induction removes every vertex on a hit path; its shortest version
has no repeated vertex. A nonempty finite closed set has a bottom strong
 component, which remains closed in the full graph. Conversely a failed sink
 component survives every removal. The standard strong-component framework is
 recorded in \cite[Theorems 3.10 and 5.2]{Tarjan}; the finite closure argument
 above is proved directly for the present arithmetic graph.
\end{proof}
No claimed efficient factorization or prime distribution theorem is involved:
a complete sieve, factorization of each $A_q<p$, and finite square-divisor boxes
terminate. There are fewer than $2p^2$ divisor candidates, as a crude bound.

\newpage
\begin{theorem}[An actual closed failure component]
At $p=2521$, the complete graph contains the sink seven-cycle
\[
 \boxed{11\to211\to683\to89\to17\to643\to113\to11.}
 \tag{2}
\]
Every vertex in this component has exactly the displayed successor, and both
of its complete local selectors fail.
\end{theorem}
\begin{proof}
All displayed numbers are prime by complete trial division. The source data are
\[
\begin{array}{r|r|l|r|r}
 q&R_q&A_q&\text{next nonresidue prime}&\tau(A_q^2)\\\hline
11&11&633=3\cdot211&211&9\\
211&211&683&683&3\\
683&683&801=3^2\cdot89&89&15\\
89&267&697=17\cdot41&17&9\\
17&51&643&643&3\\
643&643&791=7\cdot113&113&9\\
113&339&715=5\cdot11\cdot13&11&27.
\end{array}                                                \tag{3}
\]
The three targets $(-1,-p,-p^{-1})$, in this order of vertices, are
\[
 (10,9,5),\ (210,11,96),\ (682,211,560),\ (266,149,224),
 \ (50,29,44),\ (642,51,290),\ (338,191,71).
\]
The complete products of the original centered factor intervals have supports
\begin{align*}
W_{11}&=\{1,2,3,4,6,7,8\},&W_{211}&=\{1,38,50\},\\
W_{683}&=\{1,3,9,76,89,110,118,228,238,267,307,330,485,492,617\},\\
W_{89}&=\{1,17,41,46,110,163,172,238,254\},&W_{17}&=\{1,28,31\},\\
W_{643}&=\{1,7,92,108,113,126,148,239,387\},\\
W_{113}&=\{1,5,11,13,19,32,37,53,55,65,68,70,142,143,160,185,206,\\
&\hspace{22mm}209,214,232,247,265,266,275,313\}.
\end{align*}
Every target is absent from its support. The factorizations and exact character
tests give no other outgoing nonresidue prime. Thus the component has no exit.
The 75 original divisor vectors are not replaced by their 71 distinct support
positions; their full integer coefficients are retained in the certificate.
\end{proof}

The effective indices $|\langle-1,2,t:t\mid A_q\rangle/\operatorname{Stab}(W_q)|$
are respectively
\[
10,210,682,176,32,642,224;
\]
all seven stabilizers are trivial. The complete graph has 187 vertices, four
occupied vertices, and 107 trapped vertices. Its unique sink is (2).
Thus neither arbitrary factor choice nor starting at the least nonresidue
prime forces reachability of a hit.

This is not an ES counterexample:
\[
 \frac4{2521}=\frac1{636}+\frac1{5611746}+\frac1{70588}
\]
is the original M state $R=23,u=8,(h,r,s,\lambda)=(2,2,159,7)$.
A same-prime source expansion also repairs (2) explicitly. For $q=11$, allow
the separately marked quadratic-residue multiplier 5 instead of 1:
\[
 R=55,\ A=644=2^2\cdot7\cdot23,\ u=16,
 \quad(h,r,s,\lambda)=(1,4,161,3),
\]
\[
 \boxed{\frac4{2521}=\frac1{644}+\frac1{1217643}+\frac1{30252}.}
\]
The old and new $A$ values are coprime. These are new available divisor words,
not a relabelling of an old absent coefficient. No universal assertion about
this enlarged multiplier family is made.

 A primitive index-six vertex can also feed a closed failed component:
\[
 p=3361:\quad31\to53\to11\to281\to1051\to1103\to31,
\]
 with effective indices $6,104,10,560,1050,1102$. The full source includes the
 known saturated inactive factor $2^4$ at the first vertex. Growth of defect
 index is not a well-founded escape mechanism. The index-six normal-form
 antecedent is Bryan's exact combined-target analysis \cite[\S\S3--8]{BryanSix};
 the closed-component calculation here retains the complete successor graph and
 local selector data rather than attributing those new finite checks to that source.

\newpage
\paragraph{The full composite transition.}
For every edge retain $m=v_r(A_q)$, $B=A_q/r^m$, and $c=r^{m-1}B$.
Then
\[
 p+\sigma_q q=4cr,\qquad
 \boxed{4A_r=(4c+\sigma_r)r-\sigma_q q.}                  \tag{4}
\]
The successor inverse-exterior target projects to $\sigma_q q\pmod r$.
If $r\equiv1\pmod4$, the full target modulo $3r$ is the CRT pair
\[
 \boxed{-p\longleftrightarrow(2,\sigma_q q\bmod r).}       \tag{5}
\]
Its inverse is $ar(r^{-1}\bmod3)+3b(3^{-1}\bmod r)$ for pair $(a,b)$.
When $\sigma_q=3$, the integer $3q$ is a nonunit modulo $3r$:
one must retain the correction $3q-4cr$, rather than cancel three.
Reciprocity on every original prime factor gives $(t/R_q)=(t/p)$.
For an outgoing factor,
\[
 (r/q)=-1\ (q=3\pmod4),\qquad
 (r/q)=-(-3/r)\ (q=1\pmod4).
\]
These imply no two-cycle of two $3\pmod4$ primes, but do not exclude (2).

Writing $\Delta=A_r-A_q$, the exact common word domain is
\[
 \operatorname{Div}(A_q^2)\cap\operatorname{Div}(A_r^2)
 =\operatorname{Div}(\gcd(A_q,A_r)^2).
\]
On it, $\beta_t^{(r)}=\beta_t^{(q)}+v_t(A_q)-v_t(A_r)$,
$g_E'=g_E$ and $g_M'=g_M+\Delta$. The gate modulus changes as well.
No homomorphism between the two unit groups or their period quotients follows
from (4). In particular, the true M hit $u=1$ on $19\to13$ at $p=1801$
remains an available divisor but loses its gate when $A$ changes from 455 to 460.

\begin{theorem}[Exact cycle coordinates]
For a simple cycle $q_0,\ldots,q_{l-1}$ with
$p+\sigma_iq_i=4c_iq_{i+1}$, put $D=4^l\prod c_i$ and $S=\prod\sigma_i$.
For each starting index $i$, compose the affine equations, obtaining
$Dq_i=Sq_i+pT_i$. Then
\[
 \boxed{h=\gcd_iT_i,\qquad p=(D-S)/h,\qquad q_i=T_i/h.}   \tag{6}
\]
 An ordered coefficient word therefore encodes at most one such prime cycle.
\end{theorem}
\begin{proof}
Compute a prefix by $D_0=S_0=1,T_0=0$ and
$D_{j+1}=4c_{i+j}D_j$, $S_{j+1}=\sigma_{i+j}S_j$,
$T_{j+1}=\sigma_{i+j}T_j+D_j$. Induction gives the stated affine composition.
Since $p\ne q_i$ are prime, $p\mid D-S$; also $D>S$. Put $h_0=(D-S)/p$.
Then $T_i=h_0q_i$. The distinct cycle primes have common gcd one, so $h=h_0$.
Conversely, (6) gives candidates; primality, the original ranges, factor edges,
types and complete local gates must all be verified. This is a finite inverse,
 not an assertion that arbitrary coefficient words pass those tests.
 \end{proof}
 This affine cycle inverse is the all-type continuation of the primitive sextic
 recurrence in \cite[\S\S2--4]{BryanChain}; the coefficients, vertex types and
 converse checks above are stated explicitly because the presentations are not
 identified without those maps.
For (2), $D=8252375040$, $S=27$, and $h=3273453$, so $D-S=2521h$ exactly.
The contraction of the affine updates is consistent with the failed cycle.

\paragraph{Finite scope and the next implication.}
Two standard-library implementations agree on every graph for the 137 hard
primes through 50,000: 188,255 vertices, 191,789 distinct edges, 8,357,211
original divisor candidates, and 1,709 occupied vertices. Thirty-nine primes
have closed failure traps. The least hard prime with such a trap is 2521.
The code also exhausts 38,636 small directed-graph/hit-mask fixtures. These
are finite checks, not a uniform prime-coverage theorem or independent review.
The proposed universal local closure is false. A sufficient global target
would instead prove that not \emph{every} vertex is trapped, or supply a proved
escape through an explicitly enlarged original source. Neither implication
is inserted as a hypothesis here.

\begingroup\footnotesize
\begin{thebibliography}{6}
\bibitem{Graph} Bryan, C. (2026, August 15). \emph{External-nonresidue factor
cycle}. \texttt{chasebryan/centl}, \texttt{EXTERNAL-NR-FACTOR-CYCLE.md}, \S\S1--8;
commit \texttt{42508684c4386b575d7b52e763565345a76bcb9e}. Basic graph antecedent.
\bibitem{ET} Elsholtz, C., \& Tao, T. (2013). Counting the number of solutions
to the Erd\H{o}s--Straus equation on unit fractions. \emph{J. Aust. Math. Soc.,94}(1),
50--105. arXiv:1107.1010, \S2, Propositions 2.2 and 2.6.
\bibitem{BryanSix} Bryan, C. (2026, August 15). \emph{Exact index-six normal form
for the combined FAB targets}. \texttt{chasebryan/centl},
\texttt{FAB-INDEX6-COMBINED-DEFECT.md}, \S\S3--8; commit
\texttt{42508684c4386b575d7b52e763565345a76bcb9e}, blob
\texttt{81f984e6961c2b6cf4157bf8089305a5b2dbb59e}.
\bibitem{BryanChain} Bryan, C. (2026, August 15). \emph{Primitive sextic defect
chains and the neighbor-square recurrence}. Same repository and commit,
\texttt{ES-PRIMITIVE-SEXTIC-CHAIN.md}, \S\S2--4; blob
\texttt{523f133a65771fdbd5a7e5cfecc7418798dd8c35}.
\bibitem{Tarjan} Tarjan, R. E., \& Zwick, U. (2022). \emph{Finding strong
components using depth-first search}. arXiv:2201.07197v3, Theorems 3.10 and 5.2.
The original algorithm appeared in \emph{SIAM J. Comput. 1}(2), 146--160 (1972),
doi:10.1137/0201010.
\bibitem{Prior} The Clankers (2026). \emph{Turn 1: joint E/M failure structure}.
Complete predecessor proofs and hashes accompany the workbench. Its effective
index-six classification is distinguished from the earlier ambient-index result.
\end{thebibliography}
\endgroup
\end{document}
